\documentclass[11pt]{article}

% Packages
\usepackage[margin=1in]{geometry}
\usepackage{setspace}
\usepackage{parskip}
\usepackage{times}
\usepackage{fancyhdr}

% Header setup
\pagestyle{fancy}
\fancyhf{}
\lhead{}
\chead{\small Scott Nguyen --- Statement of Purpose --- Engineering Science (Mechanical \& Aerospace), UCSD/SDSU Joint Doctoral Program}
\rhead{}
\renewcommand{\headrulewidth}{0pt}

% No indent
\setlength{\parindent}{0pt}
\onehalfspacing

\begin{document}
	
	My interest in Guidance, Navigation, and Control (GNC) began during my undergraduate aerospace engineering coursework, when I first encountered orbital mechanics and how a spacecraft could determine its state autonomously. The structure of nonlinear dynamics and estimation fit naturally with how I think, and I pursued advanced courses and independent projects to deepen these skills. I realized that my strengths in analytical reasoning, coding, and modeling aligned with demands of autonomous spacecraft navigation.
	
	My academic preparation spans two graduate degrees focused on building a strong foundation for autonomous GNC research. I completed an M.S. in Aerospace Engineering at the University of Illinois Urbana–Champaign, focusing on estimation and spacecraft dynamics. I am now completing a second M.S. in Electrical/Space Systems Engineering at the University of New Mexico to strengthen my background in nonlinear filtering and sensor modeling. Across both programs, I gained experience with the Extended Kalman Filter, Unscented Kalman Filter, Monte Carlo analysis, and high-fidelity dynamic models incorporating perturbations and relative orbital elements. These efforts prepared me for advanced work in multi-sensor estimation and nonlinear guidance.
	
	My professional experiences reinforced this preparation. At the Space Dynamics Laboratory, I contributed to angles-only orbit estimation research and learned the challenges of filtering with limited observability. At Blue Canyon Technologies, I tested star trackers, inertial measurement units, reaction wheels, and solar array drive assemblies and supported flight-software verification using COSMOS and C-based S-functions. At Varda Space Industries, I implemented an EKF for capsule state estimation, developed reentry Monte Carlo tools, and evaluated gravity models for robustness. At Blue Origin, I designed hybrid Active Disturbance Rejection Control–Sliding Mode Control architectures and built closed-loop simulations for control-law validation. Most recently, at MIT Lincoln Laboratory, I have worked on radio-frequency–based state estimation, probabilistic detection methods, and S-band architectures. Together, these roles deepened my interest in autonomy and strengthened my preparation for a Ph.D.
	
	I am applying to the UCSD–SDSU Joint Doctoral Program in Engineering Science (Mechanical and Aerospace) to complete a Ph.D. focused on autonomous navigation, multi-sensor fusion, and trajectory design for cislunar and deep-space missions. My goal is to develop navigation architectures and simulation frameworks that allow spacecraft to operate independently of the Global Positioning System and the Deep Space Network. I hope to work as a GNC engineer advancing autonomous flight systems.
	
	My research interests include image-based navigation, terrain-relative feature detection, multi-sensor state estimation, nonlinear filtering, trajectory design in the Circular Restricted Three-Body Problem, and uncertainty quantification. These areas align closely with Dr. Pablo Machuca’s work, and the program’s structure—UCSD’s theoretical depth combined with SDSU’s applied engineering environment—matches how I approach research in mechanical and aerospace systems.
	
	Growing up in a low-income immigrant household taught me independence, resilience, and resourcefulness. These qualities, combined with my technical preparation, show that I am prepared to contribute to the UCSD–SDSU Joint Doctoral Program.
	
\end{document}